%% ========================================================================
%% Appendix A: Referensi Perintah Linux
%% ========================================================================

\chapter{Referensi Perintah Linux}
\label{app:command-reference}

Lampiran ini merangkum perintah Linux yang paling sering digunakan sepanjang buku.
Gunakan lampiran ini sebagai panduan cepat, lalu baca \cmd{man nama-perintah}
atau \cmd{nama-perintah --help} ketika kamu memerlukan penjelasan yang lebih rinci.

%% ------------------------------------------------------------------------
\section{Perintah Sistem File}

\paragraph*{ls --- Menampilkan Isi Direktori}
\cmd{ls} menampilkan isi direktori dalam berbagai format.
\begin{lstlisting}[language=bash]
ls -la
ls -lh
ls -lt
\end{lstlisting}

\paragraph*{cd --- Berpindah Direktori}
\cmd{cd} digunakan untuk berpindah dari satu direktori ke direktori lain.
\begin{lstlisting}[language=bash]
cd /etc
cd ..
cd ~
cd -
\end{lstlisting}

\paragraph*{pwd --- Menampilkan Direktori Saat Ini}
\cmd{pwd} menampilkan path direktori kerja saat ini.
\begin{lstlisting}[language=bash]
pwd
\end{lstlisting}

\paragraph*{tree --- Menampilkan Struktur Direktori}
\cmd{tree} menampilkan isi direktori dalam bentuk pohon.
\begin{lstlisting}[language=bash]
tree
tree -L 2
tree ~/lab-os
\end{lstlisting}

\paragraph*{stat --- Menampilkan Metadata Berkas}
\cmd{stat} menampilkan ukuran, inode, permission, dan waktu berkas.
\begin{lstlisting}[language=bash]
stat laporan.txt
stat /etc/passwd
\end{lstlisting}

%% ------------------------------------------------------------------------
\section{Perintah Manajemen File}

\paragraph*{cp --- Menyalin File dan Direktori}
\cmd{cp} digunakan untuk menyalin file atau direktori.
\begin{lstlisting}[language=bash]
cp file1.txt salinan.txt
cp laporan.txt ~/arsip/
cp -r proyek/ backup-proyek/
\end{lstlisting}

\paragraph*{mv --- Memindahkan dan Mengganti Nama}
\cmd{mv} memindahkan file atau mengganti namanya.
\begin{lstlisting}[language=bash]
mv draft.txt final.txt
mv laporan.txt ~/arsip/
mv gambar/ ~/backup/
\end{lstlisting}

\paragraph*{rm --- Menghapus File dan Direktori}
\cmd{rm} menghapus file, dan dengan opsi tertentu juga direktori.
\begin{lstlisting}[language=bash]
rm catatan.txt
rm -r folder-lama/
rm -rf cache-sementara/
\end{lstlisting}

\paragraph*{mkdir --- Membuat Direktori}
\cmd{mkdir} membuat satu atau lebih direktori baru.
\begin{lstlisting}[language=bash]
mkdir arsip
mkdir -p proyek/{dokumen,skrip,output}
mkdir -p ~/lab-os/chapter04-files
\end{lstlisting}

\paragraph*{touch --- Membuat Berkas Kosong dan Mengubah Timestamp}
\cmd{touch} membuat file kosong atau memperbarui timestamp file yang ada.
\begin{lstlisting}[language=bash]
touch catatan.txt
touch file1.txt file2.txt
touch -t 202605091030 laporan.txt
\end{lstlisting}

\paragraph*{ln --- Membuat Hard Link dan Symbolic Link}
\cmd{ln} membuat hard link atau symbolic link.
\begin{lstlisting}[language=bash]
ln sumber.txt salinan-hard.txt
ln -s /var/log/syslog log-sistem
ln -s ~/lab-os/chapter04-files proyek-link
\end{lstlisting}

\paragraph*{find --- Mencari File}
\cmd{find} mencari file berdasarkan nama, tipe, ukuran, atau waktu.
\begin{lstlisting}[language=bash]
find . -name "*.txt"
find ~/lab-os -type d
find /var/log -mtime -1
\end{lstlisting}

\paragraph*{locate --- Mencari File dari Basis Data}
\cmd{locate} mencari file dengan cepat dari basis data yang sudah diindeks.
\begin{lstlisting}[language=bash]
locate bashrc
locate nginx.conf
sudo updatedb
\end{lstlisting}

\paragraph*{tar --- Membuat dan Mengekstrak Arsip}
\cmd{tar} dipakai untuk membuat, melihat, dan mengekstrak arsip.
\begin{lstlisting}[language=bash]
tar -czf backup.tar.gz data-sumber/
tar -xzf backup.tar.gz
tar -tf backup.tar.gz
\end{lstlisting}

\paragraph*{gzip --- Mengompresi File}
\cmd{gzip} mengompresi file tunggal menggunakan format gzip.
\begin{lstlisting}[language=bash]
gzip laporan.log
gzip -k laporan.log
gunzip laporan.log.gz
\end{lstlisting}

\paragraph*{zip --- Membuat Arsip ZIP}
\cmd{zip} membuat arsip dalam format ZIP.
\begin{lstlisting}[language=bash]
zip arsip.zip laporan.txt
zip -r proyek.zip proyek/
zip -9 hasil.zip data.csv
\end{lstlisting}

\paragraph*{unzip --- Mengekstrak Arsip ZIP}
\cmd{unzip} mengekstrak isi arsip ZIP.
\begin{lstlisting}[language=bash]
unzip arsip.zip
unzip arsip.zip -d hasil-ekstrak/
unzip -l arsip.zip
\end{lstlisting}

%% ------------------------------------------------------------------------
\section{Perintah Manajemen Proses}

\paragraph*{ps --- Menampilkan Proses}
\cmd{ps} menampilkan snapshot proses yang sedang berjalan.
\begin{lstlisting}[language=bash]
ps
ps aux
ps aux | grep ssh
\end{lstlisting}

\paragraph*{top --- Memantau Proses Secara Real-Time}
\cmd{top} memantau penggunaan CPU, memori, dan proses aktif.
\begin{lstlisting}[language=bash]
top
top -b -n 1
top -u namauser
\end{lstlisting}

\paragraph*{htop --- Pemantauan Proses Interaktif}
\cmd{htop} menyediakan tampilan proses yang lebih interaktif dari \cmd{top}.
\begin{lstlisting}[language=bash]
htop
htop -u namauser
\end{lstlisting}

\paragraph*{kill --- Mengirim Sinyal ke Proses}
\cmd{kill} mengirim sinyal ke proses berdasarkan PID.
\begin{lstlisting}[language=bash]
kill 1234
kill -TERM 1234
kill -9 1234
\end{lstlisting}

\paragraph*{pkill --- Menghentikan Proses Berdasarkan Nama}
\cmd{pkill} mengirim sinyal ke proses berdasarkan nama proses.
\begin{lstlisting}[language=bash]
pkill firefox
pkill -f "python3 -m http.server"
pkill -u namauser bash
\end{lstlisting}

\paragraph*{jobs --- Menampilkan Job Shell}
\cmd{jobs} menampilkan job yang sedang dikelola oleh shell saat ini.
\begin{lstlisting}[language=bash]
jobs
jobs -l
\end{lstlisting}

\paragraph*{bg --- Menjalankan Job di Latar Belakang}
\cmd{bg} melanjutkan job yang berhenti ke mode background.
\begin{lstlisting}[language=bash]
bg
bg %1
\end{lstlisting}

\paragraph*{fg --- Membawa Job ke Foreground}
\cmd{fg} memindahkan job background ke foreground.
\begin{lstlisting}[language=bash]
fg
fg %1
\end{lstlisting}

\paragraph*{nice --- Menjalankan Proses dengan Prioritas Tertentu}
\cmd{nice} menjalankan proses dengan nilai prioritas CPU tertentu.
\begin{lstlisting}[language=bash]
nice -n 10 ./proses-berat.sh
nice -n -5 python3 analisis.py
\end{lstlisting}

\paragraph*{renice --- Mengubah Prioritas Proses}
\cmd{renice} mengubah nilai prioritas proses yang sudah berjalan.
\begin{lstlisting}[language=bash]
renice 5 -p 1234
sudo renice -10 -p 1234
\end{lstlisting}

%% ------------------------------------------------------------------------
\section{Perintah Jaringan}

\paragraph*{ip --- Menampilkan dan Mengatur Jaringan}
\cmd{ip} digunakan untuk melihat interface, alamat, dan route jaringan.
\begin{lstlisting}[language=bash]
ip addr show
ip link show
ip route
\end{lstlisting}

\paragraph*{ss --- Menampilkan Socket dan Port}
\cmd{ss} menampilkan koneksi jaringan dan port yang aktif.
\begin{lstlisting}[language=bash]
ss -tlnp
ss -an
ss -tulpn | grep :22
\end{lstlisting}

\paragraph*{ping --- Menguji Konektivitas}
\cmd{ping} menguji apakah host lain dapat dijangkau melalui jaringan.
\begin{lstlisting}[language=bash]
ping 8.8.8.8
ping -c 4 google.com
ping -c 2 localhost
\end{lstlisting}

\paragraph*{curl --- Mengambil Data dari URL}
\cmd{curl} mengambil data dari URL dan sering dipakai untuk menguji API.
\begin{lstlisting}[language=bash]
curl https://example.com
curl -I https://example.com
curl -O https://example.com/file.txt
\end{lstlisting}

\paragraph*{wget --- Mengunduh File}
\cmd{wget} digunakan untuk mengunduh file dari web atau FTP.
\begin{lstlisting}[language=bash]
wget https://example.com/file.iso
wget -c https://example.com/file.iso
wget -O salinan.html https://example.com
\end{lstlisting}

\paragraph*{dig --- Memeriksa DNS}
\cmd{dig} memeriksa resolusi DNS dan catatan domain.
\begin{lstlisting}[language=bash]
dig example.com
dig example.com MX
dig @8.8.8.8 example.com
\end{lstlisting}

\paragraph*{scp --- Menyalin File lewat SSH}
\cmd{scp} menyalin file antarhost melalui koneksi SSH.
\begin{lstlisting}[language=bash]
scp file.txt user@server:/home/user/
scp user@server:/var/log/syslog .
scp -r proyek/ user@server:/home/user/
\end{lstlisting}

\paragraph*{ssh --- Mengakses Host Jarak Jauh}
\cmd{ssh} membuka sesi shell aman ke host lain.
\begin{lstlisting}[language=bash]
ssh user@server
ssh -p 2222 user@server
ssh -i ~/.ssh/id_ed25519 user@server
\end{lstlisting}

%% ------------------------------------------------------------------------
\section{Perintah Sistem}

\paragraph*{uname --- Menampilkan Informasi Kernel}
\cmd{uname} menampilkan informasi kernel dan arsitektur sistem.
\begin{lstlisting}[language=bash]
uname -a
uname -r
uname -m
\end{lstlisting}

\paragraph*{df --- Menampilkan Kapasitas Filesystem}
\cmd{df} menampilkan kapasitas dan penggunaan filesystem yang ter-mount.
\begin{lstlisting}[language=bash]
df -h
df -T
df -h /home
\end{lstlisting}

\paragraph*{du --- Menampilkan Ukuran Direktori}
\cmd{du} menghitung ukuran file atau direktori.
\begin{lstlisting}[language=bash]
du -sh ~/lab-os
du -h --max-depth=1 /var/log
du -sh *
\end{lstlisting}

\paragraph*{free --- Menampilkan Penggunaan Memori}
\cmd{free} menampilkan penggunaan RAM dan swap.
\begin{lstlisting}[language=bash]
free -h
free -m
watch -n 1 free -h
\end{lstlisting}

\paragraph*{uptime --- Menampilkan Lama Sistem Menyala}
\cmd{uptime} menampilkan lama sistem aktif dan beban rata-rata.
\begin{lstlisting}[language=bash]
uptime
uptime -p
uptime -s
\end{lstlisting}

\paragraph*{journalctl --- Membaca Log systemd}
\cmd{journalctl} membaca log yang dikelola oleh systemd journal.
\begin{lstlisting}[language=bash]
journalctl -b
journalctl -u ssh
journalctl -u nginx -n 50
\end{lstlisting}

\paragraph*{systemctl --- Mengelola Layanan}
\cmd{systemctl} mengelola layanan dan unit systemd.
\begin{lstlisting}[language=bash]
systemctl status ssh
sudo systemctl restart nginx
systemctl list-units --type=service
\end{lstlisting}

\paragraph*{lsblk --- Menampilkan Perangkat Blok}
\cmd{lsblk} menampilkan disk, partisi, dan mount point.
\begin{lstlisting}[language=bash]
lsblk
lsblk -f
lsblk -o NAME,SIZE,FSTYPE,MOUNTPOINT
\end{lstlisting}

\paragraph*{mount --- Me-mount Filesystem}
\cmd{mount} memasang filesystem ke sebuah mount point.
\begin{lstlisting}[language=bash]
mount | grep /dev/sdb1
sudo mount /dev/sdb1 /mnt/data
sudo mount -o ro /dev/sdb1 /mnt/data
\end{lstlisting}

\paragraph*{umount --- Melepas Mount Filesystem}
\cmd{umount} melepaskan filesystem dari mount point.
\begin{lstlisting}[language=bash]
sudo umount /mnt/data
sudo umount /dev/sdb1
\end{lstlisting}
